% method.tex - Neural Networks format

\section{Method}
\label{sec:method}

\subsection{Architectural Overview}

Censor is a biomimetic dual-pathway neural architecture implemented in PyTorch. The system processes 16-frame video clips at 224$\times$224 resolution through two complementary pathways before fusion and classification.

\textbf{Fast Pathway}: A 3D ResNet-18 variant operating on optical flow, analogous to the subcortical visual route. This pathway captures rapid motion patterns with aggressive temporal downsampling, producing a 512-dimensional feature vector.

\textbf{Slow Pathway}: A 3D Swin Transformer processing RGB frames supplemented with remote photoplethysmography (rPPG) signals, analogous to the cortical visual route. This pathway maintains high spatial resolution throughout, yielding a 768-dimensional representation.

\textbf{Fusion}: Features from both pathways are combined through squeeze-excitation attention, followed by amygdala-inspired gating that modulates the fused representation.

\textbf{Classification}: A noisy top-2 Mixture-of-Experts (MoE) mechanism routes inputs to specialized expert networks, enabling category-specific feature discrimination.

\subsection{Biomimetic Preprocessing}

\subsubsection{Optical Flow Extraction}

The fast pathway input is computed via Dual TV-L1 optical flow~\cite{wang2015offapexnet}, optimizing:
\begin{equation}
E(\vect{u}) = \int_\Omega \left( |\nabla u_1| + |\nabla u_2| \right) d\vect{x} + \lambda \int_\Omega \rho(\vect{x}, \vect{u}) d\vect{x}
\end{equation}
where $\vect{u} = (u_1, u_2)$ is the displacement field and $\rho$ is the brightness constancy term. TV regularization preserves motion discontinuities at facial muscle boundaries.

\subsubsection{rPPG Extraction}

Remote photoplethysmography extracts cardiac signals from subtle color variations:
\begin{equation}
\vect{C}(t) = 0.77 \cdot R(t) - 0.51 \cdot G(t) - 0.26 \cdot B(t)
\end{equation}
The chrominance signal is bandpass filtered (0.5--4.0 Hz) to isolate heart rate range, providing physiological arousal cues.

\subsection{Fast Pathway: 3D ResNet-18}

The fast pathway comprises three stages with increasing channel dimensions and temporal downsampling:
\begin{itemize}
\item Stage 1: Input $\dim{2\times 16 \times 112 \times 112}$, output $\dim{64 \times 8 \times 56 \times 56}$
\item Stage 2: Two residual blocks, output $\dim{128 \times 4 \times 28 \times 28}$
\item Stage 3: Two residual blocks, output $\dim{256 \times 2 \times 14 \times 14}$
\item Global pooling yields $\vect{f}_{\text{fast}} \in \dim{512}$
\end{itemize}

Residual blocks follow:
\begin{equation}
\vect{z}^{(i+1)} = \relu\left( \vect{z}^{(i)} + \mathcal{F}(\vect{z}^{(i)}; \vect{W}^{(i)}) \right)
\end{equation}
where $\mathcal{F}$ comprises 3D convolutions with batch normalization. Dropout (rate 0.1) provides regularization.

\subsection{Slow Pathway: 3D Swin Transformer}

The slow pathway processes concatenated RGB and rPPG signals ($\dim{6 \times 16 \times 224 \times 224}$) through four hierarchical stages:
\begin{itemize}
\item Stage 1: 96-dim embeddings, $(8 \times 14 \times 14)$ window resolution
\item Stage 2: 192-dim, $(8 \times 7 \times 7)$ windows after patch merging
\item Stage 3: 384-dim, $(4 \times 7 \times 7)$ windows
\item Stage 4: 768-dim, $(2 \times 7 \times 7)$ windows, output $\vect{f}_{\text{slow}} \in \dim{768}$
\end{itemize}

Shifted-window multi-head self-attention (SW-MSA) enables cross-window connections:
\begin{equation}
\text{Attention}(\vect{Q}, \vect{K}, \vect{V}) = \softmax\left( \frac{\vect{Q}\vect{K}^\top}{\sqrt{d_k}} + \vect{B} \right) \vect{V}
\end{equation}
where $\vect{B}$ encodes relative position bias.

\subsection{Fusion and Classification}

\subsubsection{Feature Fusion}

Concatenated features $\vect{f}_{\text{concat}} = [\vect{f}_{\text{fast}}; \vect{f}_{\text{slow}}] \in \dim{1280}$ are projected to $\dim{1024}$ through squeeze-excitation attention:
\begin{equation}
\vect{f}_{\text{fused}} = \sigma(\vect{W}_2 \cdot \relu(\vect{W}_1 \cdot \text{GAP}(\vect{f}_{\text{concat}}))) \cdot \vect{f}_{\text{concat}}
\end{equation}

\subsubsection{Mixture of Experts}

Noisy top-2 gating selects expert networks:
\begin{equation}
G(\vect{x}) = \softmax\left( \vect{W}_g \vect{x} + \epsilon \right), \quad \epsilon \sim \mathcal{N}(0, \sigma^2)
\end{equation}
Only the top-2 experts are activated per input, enabling specialization while maintaining efficiency.

\subsection{Training Objective}

The total loss combines classification and auxiliary terms:
\begin{equation}
\loss{total} = \crossentropy + \lambda_{\text{moe}} \loss{moe}{load} + \lambda_{\text{au}} \loss{AU}{aux}
\end{equation}
where $\loss{moe}{load}$ balances expert utilization and $\loss{AU}{aux}$ provides action unit supervision. ArcFace angular margin loss optionally constrains feature geometry:
\begin{equation}
\loss{ArcFace} = -\log \frac{e^{s \cos(\theta_{y_i} + m)}}{e^{s \cos(\theta_{y_i} + m)} + \sum_{j \neq y_i} e^{s \cos(\theta_j)}}
\end{equation}